\documentclass[letterpaper,11pt]{article}
\usepackage[T1]{fontenc}
\usepackage[utf8]{inputenc}
\usepackage{lmodern}
\usepackage[french]{babel}
\usepackage[margin=1in,headheight=28pt]{geometry}
\renewcommand{\familydefault}{\sfdefault}
\usepackage{amsmath,amssymb,amsthm,bm,enumitem}
\usepackage[dvipsnames]{xcolor}
\usepackage[small,compact]{titlesec}
\usepackage{fancyhdr,fancyvrb}
\usepackage[colorlinks=true,allcolors=blue]{hyperref}
\pagestyle{fancy}
\fancyhead[L]{STT-2100 \\ Régression}
\fancyhead[R]{A.-A. Vallée \\ Université Laval}
\fancyfoot[C]{\thepage}
\setlength{\parindent}{0pt}
\setlength{\parskip}{5pt}
\fvset{fontsize=\small,frame=single,framesep=2mm}
\newcommand{\esp}{\mathbb{E}}
\newcommand{\var}{\operatorname{Var}}
\newcommand{\bo}[1]{\bm{#1}}
\newcommand{\N}{\mathcal{N}}
% Définir \avecsolutions avant l'input pour une version corrigée.
\newcommand{\prof}[1]{\ifdefined\avecsolutions\par{\color{blue}\textbf{Solution :} #1}\par\fi}
\begin{document}
\begin{center}{\Large Exercices - Chapitre 9}\end{center}
\begin{enumerate}[leftmargin=*,wide,label=\bfseries Exercice \arabic*:]
\item Une matrice de modèle contient la constante et trois variables telles que $x_3=2x_1-x_2$. Décrivez son rang possible et la conséquence sur l'identification des coefficients. Les valeurs ajustées sont-elles nécessairement indéterminées? Comment R signale-t-il les coefficients non identifiables?
\prof{Le rang est au plus trois pour quatre colonnes. Les coefficients ne sont pas tous identifiables, mais la projection sur l'espace des colonnes et les valeurs ajustées sont uniques. \texttt{lm()} donne un ou plusieurs coefficients \texttt{NA}; \texttt{alias()} expose les dépendances.}
\item Pour la régression auxiliaire d'une variable sur les autres, $R_j^2=0{,}90$. Calculez le VIF. Si la variance de son coefficient serait $0{,}04$ avec des variables orthogonales et la même dispersion de cette variable, quelle est sa variance avec ce VIF? Quel est le facteur d'inflation de l'erreur-type?
\prof{$VIF=10$, variance $0{,}40$ et facteur d'inflation de l'erreur-type $\sqrt{10}=3{,}1623$.}
\item Pourquoi des corrélations par paires modérées ne suffisent-elles pas à exclure une multicolinéarité? Construisez un exemple avec $x_3=x_1+x_2$ et $x_1,x_2$ non corrélés de même variance.
\prof{La dépendance peut impliquer trois variables ou davantage. Dans l'exemple, $\operatorname{Corr}(x_1,x_2)=0$, tandis que les deux autres corrélations valent $1/\sqrt2$, sans corrélation parfaite par paire. Pourtant, la relation entre les trois colonnes est exacte.}
\item Les valeurs propres d'une matrice de Gram à colonnes normalisées sont $(2{,}7,0{,}299,0{,}001)$. Calculez le plus grand indice de conditionnement. Quelle distinction faites-vous entre ce nombre et le conditionnement de la matrice de Gram? Quel autre diagnostic aide à identifier les coefficients concernés?
\prof{$\eta_{max}=\sqrt{2{,}7/0{,}001}=51{,}962$. Le conditionnement spectral de la matrice de Gram vaut 2700, soit le carré. Examinez les proportions de décomposition de la variance; plusieurs coefficients ayant des proportions élevées sur une même direction mal conditionnée indiquent la dépendance concernée.}
\item Expliquez pourquoi centrer une variable avant d'inclure son carré peut aider l'interprétation et le calcul, sans garantir l'élimination de toute multicolinéarité. Comparez le retrait d'une variable, la régression ridge et la régression sur composantes principales.
\prof{Le centrage déplace la référence et peut réduire la dépendance entre termes construits, mais n'enlève pas une redondance intrinsèque entre mesures. Le retrait change la question et peut introduire un biais d'omission. Ridge conserve les variables avec une pénalité quadratique; la régression sur composantes utilise des directions orthogonales et un choix de dimension. Les deux dernières méthodes peuvent introduire un biais pour réduire la variance.}
\item Avec \texttt{mtcars}, calculez les VIF pour \texttt{wt}, \texttt{hp} et \texttt{disp} par trois régressions auxiliaires. Comparez les erreurs-types dans le modèle complet et après retrait de \texttt{disp}. Pourquoi cette comparaison ne suffit-elle pas à justifier le retrait?
\prof{Pour chaque variable, ajustez sa régression sur les deux autres puis calculez $1/(1-R_j^2)$. Le retrait peut changer les effets partiels estimés et la question scientifique; il faut examiner le contexte, les prédictions et les diagnostics.}
\end{enumerate}
\end{document}
